% page 569 du fac-simile (image p-568 du scan)
% bloc 167.6 x 263.6 mm ; marge gauche 34.4 mm
\pgc{12.700mm}{0.000mm}{
\l{}
\l{}
\l{}
\l{\cel{55}561}
\l{}
\l{\sou{sufuziono}. (\sou{medic.})\cel{2}Adfluego da sango sub la pelo. -}
\l{}
\l{\sou{sugar.(trans.})\cel{2}I.Aspirar, per la labii, la suko quan kontenas}
\l{\cel{3}ula substanci. - II. Presar, per la labii, (a) (ulo, frukto,}
\l{\cel{3}e c.) por extraktar olua suko ; (b) (\sou{aludante infanteto}) la}
\l{\cel{3}mamo di lua patrino, por extraktar lua lakto. - D.}
\l{}
\l{\sou{sugestar}. (\sou{trans.,}\cel{1}ad). Venigar en la penso-domeno. - DEFIS.}
\l{}
\l{\sou{suko}. Liquoro quan kontenas la substanco di frukti, karni, legumi,}
\l{\cel{3}e quan onu povas extraktar per preso, macero, koquo, e c.-FIR.}
\l{}
\l{\sou{sukombar}. (\sou{netrans., ad)}\cel{2}Sinkar, falar, efike da kargajo tro}
\l{\cel{3}pezoza, da koakto tro granda. - EFIS.}
\l{}
\l{\sou{sukro}. Substanco qua saporas tre dolce, extraktata de plura veje-}
\l{\cel{3}tanti (precipue kano e betravo) e qua transformesas, efike}
\l{\cel{3}da fairo, aden substanco kristalatra, solvebla en aquo. - II.}
\l{\cel{3}(\sou{kemio}).Omna substanco qua esas apta fermentacar e transform-}
\l{\cel{3}esar aden alkoholo e CO2. - DEFIRS.}
\l{}
\l{\sou{sukubo}. (\sou{teol.}) Demono noktala qua igas su aspektar quale muliero}
\l{\cel{3}por koitar kun viro. - DEFS.}
\l{}
\l{\sou{sukusar}. (\sou{trans.)}\cel{2}I. Agitar bruske en omna lua parti. - II.}
\l{\cel{3}(\sou{metaf}.)Agitar bruske (kozo) por desintrikar su de olu. - FI.}
\l{}
\l{\sou{sulo}. Extensajo sur qua su apogas la korpi, surface di Tero.- EFIS}
\l{}
\l{\sou{sulfo}. (\sou{kemio}). Korpo simpla (metaloida), solida, sen-sapora,}
\l{\cel{3}klare-flava, tre inflamebla, e qua, kande kombustata, exhalas}
\l{\cel{3}odoro tre forta e sufokiva. -\cel{1}\sou{Simbolo kemiala}\cel{1}:S.\cel{2}- EFIS.}
\l{}
\l{\sou{sulko}. (\sou{agrokultivo)}.\cel{2}Trancheo quan facas la soko di plugilo. IS}
\l{}
\l{\sou{sultano}. (\sou{olim)}\cel{2}La suvereno di Turkia. - DEFIRS.}
\l{}
\l{\sou{sumo}. Quanto ek du o plura quanti quin onu adicionis. - DEFIRS.}
\l{}
\l{\sou{sumako}. (\sou{bot.}) Planto ek la familio di la terebinto, uzata da la}
\l{\cel{3}mejisisti. - L.\cel{1}\sou{rhus.}\cel{2}- DEFIRS.}
\l{}
\l{\sou{sumaria}. Rapid-agema ; sen disipo di tempo e sen formalaji ne-}
\l{\cel{3}utila. - EF.}
\l{}
\l{\sou{sumnar.(trans, pri)}. (\sou{yurocienco}). Imperar (ad ulu) ke lu agez to}
\l{\cel{3}quon onu havas la yuro postular de lu, minacante lu ye la}
\l{\cel{3}puniso a qua lu su expozus se lu refuzus. - EF.}
\l{}
\l{\sou{suno}. I.Astro lumifanta, qua esas la centro di la planetaro di}
\l{\cel{3}qua Tero esas parto. - II. (\sou{metaf.}) To quo brilas quale Suno.-}
\l{\cel{3}DE.}
\l{}
\l{sundio.\cel{2}La dio unesma di la semano - konsakrita, da la eklezio}
\l{\cel{3}kristana, a repozo ed a la kulto di la deo kristana, memorige}
\l{\cel{3}di la rezurekto di Iesu Kristo. - DE.}
}
